\chapter{Allen-Cahn energy and minimal interfaces}\label{sec:allen-cahn}
Allen--Cahn theory approaches geometric interfaces through functions. Chodosh's 2019 course notes, in the supplied revision dated June 29, 2023, give the basic analytic and variational setup \cite{ChAC}. The energy below is a diffuse-interface energy; its min-max levels are not identified with cycle-space widths without a comparison theorem.

\begin{definition}[The Allen-Cahn energy]\label{def:allen-cahn-energy}
Let $(M^{n+1},g)$ be a smooth closed manifold and $\varepsilon>0$. For $u\in H^1(M)\cap L^4(M)$ set
\[
 E_\varepsilon(u)=\int_M\left(\frac\varepsilon2|\nabla u|^2+\frac{W(u)}\varepsilon\right)dV_g,
 \qquad W(s)=\frac14(1-s^2)^2.
\]
We assign energy $+\infty$ outside this domain. The wells are $u=\pm1$. The interface normalization is
$c_W=\int_{-1}^{1}\sqrt{2W(s)}\,ds=2\sqrt2/3$.
Source: \cite{ChAC}, Definition 2.1 and Section 2.2. Any comparison with area must keep this normalization.
\end{definition}

\begin{proposition}[Equation and second variation of Allen-Cahn]\label{prop:allen-cahn-variation}
A smooth function is critical for $E_\varepsilon$ if and only if
$\varepsilon^2\Delta_g u=W'(u)=u(u^2-1)$.
At a smooth critical point its second variation in the direction $\phi$ is
\[
 Q_\varepsilon(\phi,\phi)=\int_M\left(\varepsilon|\nabla\phi|^2+\varepsilon^{-1}W''(u)\phi^2\right)dV_g.
\]
Compare \cite{ChAC}, Section 2.1.
\end{proposition}
\begin{proof}
Differentiate the integrand at $u+t\phi$ once and integrate the gradient term by parts on the closed manifold. The first derivative is
$\int_M(-\varepsilon\Delta_g u+\varepsilon^{-1}W'(u))\phi$.
Its vanishing for every smooth test function is exactly the equation. Differentiating a second time gives the displayed quadratic form.
\end{proof}

\begin{example}[A single transition layer]\label{ex:allen-cahn-layer}
For the potential in Definition \ref{def:allen-cahn-energy},
$q(t)=\tanh(t/\sqrt2)$ satisfies $q''=W'(q)$, $q(\pm\infty)=\pm1$, and $q'=\sqrt{2W(q)}$.
Thus $u(t)=q(t/\varepsilon)$ solves the one-dimensional equation, and its energy per unit transverse area is
\[
 \int_{\mathbb R}\left(\tfrac12|q'|^2+W(q)\right)dt
 =\int_{-1}^{1}\sqrt{2W(s)}\,ds=c_W.
\]
The substitution $s=q(t)$ proves the identity and explains why diffuse-interface energies must be divided by $c_W$ when comparing with area. See \cite{ChAC}, Section 2.2.
\end{example}

\begin{remark}[Multiplicity, curvature, and index in dimension three]\label{rem:allen-cahn-minmax}
The supplied Chodosh--Mantoulidis paper \cite{CM} is arXiv:1803.02716v2, with a manuscript date of May 16, 2018. Its Sections 3--5 develop curvature estimates, separation of transition sheets, and the relationship between multiplicity and Jacobi fields for solutions with controlled energy and Morse index. Theorem 1.7 explains why a smooth limiting component with multiplicity, or a one-sided component, forces a positive Jacobi field (on the two-sided cover when needed). For bumpy metrics on closed three-manifolds this obstruction gives two-sided multiplicity-one components. The full convergence and index theorems, including their analytic hypotheses, remain external dependencies here. We do not infer any equality of Allen--Cahn and Almgren--Pitts widths from this source discussion.
\end{remark}
